\section{Evaluation}
\label{sec:evaluation}
We evaluate three questions: how effectively the analysis detects type
errors (RQ1), how its mechanisms affect decisions and cost (RQ2), and what
new errors it identifies in independent repositories (RQ3).

\subsection{Experimental Setup}
\label{sec:setup}
\paragraph{Benchmark.}
The historical benchmark contains xxx public GitHub bug-fix records from
xxx Python repositories, collected over xxx. Each record identifies the
buggy revision, its repaired revision, the target operation or annotation,
and the evidence linking the change to the error. Table~\ref{tab:dataset}
summarizes the corpus. The detector analyzes each revision separately using that revision's
source and available dependencies. Patch descriptions, cross-revision
comparisons, and issue metadata are reserved for assessment. The corpus emphasizes recent public type errors. We assign
runtime and annotation-contract labels separately, retaining a link when
a record exhibits both categories.

A target's category and entry scope determine its eligibility. Runtime
targets satisfy the protocol-failure definition in Section~\ref{sec:background};
annotation targets concern the semantic compatibility of a declaration
with behavior. For every tool, results retain the full benchmark denominator
and the applicable-subset denominator. Unresolved targets, timeouts, and
unsupported cases have separate counts.

\begin{table}[t]
\centering
\caption{Benchmark composition. Category counts can overlap; unique records are counted once.}
\label{tab:dataset}
\begin{tabular}{lr}
\toprule
Characteristic & Count or value\\
\midrule
Historical repositories & xxx\\
Unique bug-fix records & xxx\\
Runtime type-error records & xxx\\
Annotation-contract error records & xxx\\
Records with both categories & xxx\\
Collection interval & xxx\\
Independent discovery repositories & xxx\\
\bottomrule
\end{tabular}
\end{table}

\paragraph{Baselines.}
We compare Mypy, Pyright~\cite{pyright}, Pyrefly~\cite{pyrefly}, Pyinder,
MOPSA on its supported subset,
direct LLM detection, and annotation inference followed by a type checker.
The checker baselines represent declared and inferred type checking;
Pyinder addresses the same detection task~\cite{oh2024pyinder}; MOPSA
provides a comparison with relational Python analysis~\cite{monat2020python}.
The annotation pipeline uses xxx followed by xxx and evaluates errors in
the original program separately from inconsistencies introduced by generated
annotations. Direct LLM detection uses xxx with the entry and repository
context specified by xxx. Tool versions and configurations are xxx.
Mypy uses \texttt{--check-untyped-defs} for checking unannotated function
bodies~\cite{mypyCommandLine}.

\paragraph{Configuration.}
The analysis uses LLM xxx, decoding configuration xxx, and prompts xxx.
Python and dependency versions are xxx, pyflow revision xxx, and hardware
xxx. Entry selection and dependency models are xxx. The limits for explicit
paths, implicit type alternatives, loop/recursive unfolding, conflict
refinement, and wall-clock time are xxx. The aggregation rule for completion, resolution, and outcome
counts across runs is xxx. Cost accounting includes all runs. These settings apply to the comparisons below.

\paragraph{Detection and precision.}
Paired detection requires a target report in the buggy revision and removal
of the corresponding conflict in the fixed revision, with both analyses
complete and the fixed target semantically resolved. For runtime errors,
the buggy outcome must satisfy the entry and exception-escape criteria.
For annotation-contract errors, the repaired declaration or behavior must
resolve the confirmed incompatibility. A fixed target that becomes
unreachable requires reachability evidence. Let $N_k$ be the number of
eligible records for category $k$ and $D_k$ its paired detections. We report
$D_k/N_k$ together with the full-pool rate. A record in both categories
contributes once to each category and once to any unique-record aggregate.

Diagnostic precision is the fraction of reviewed reports confirmed as
actual errors, including non-target errors. We report it separately for
the two categories. Review follows the source, declared entry, repair or
other behavioral evidence, and exception disposition; the adjudication
procedure and reviewer agreement are xxx. Duplicate reports of the same
root cause and category are grouped within a revision. Cross-category
links are retained rather than counted as duplicate diagnoses.

\paragraph{Resolution and cost.}
Completion means the analysis process finishes within its resource limits;
resolution means it establishes a target's semantic disposition. We report
both, together with unresolved, timed-out, and unsupported targets.
Total cost includes call-graph construction, fact and relation extraction,
composition, LLM requests, trace expansion, re-inference, and caller
reanalysis. Measurements include wall-clock time, peak memory, requests,
tokens, summary applications, and expanded source size. Medians and tail
latencies are computed over the explicitly identified completed set;
timeouts retain their consumed budget in total-cost accounting.

\subsection{RQ1: Detection Effectiveness}
\label{sec:rq1}
Table~\ref{tab:detection} reports paired detection and diagnostic precision
by category. Our approach detects xxx runtime targets and xxx annotation
contract targets, with precision of xxx\% and xxx\%, respectively. The
corresponding full-pool paired rates are xxx\% and xxx\%.

\begin{table}[t]
\centering
\small
\setlength{\tabcolsep}{4pt}
\caption{Detection effectiveness. Paired counts and precision (\%) are separated by error category. MOPSA uses its supported subset of xxx records, covering xxx\% of the full pool.}
\label{tab:detection}
\begin{tabular}{lrrrr}
\toprule
& \multicolumn{2}{c}{Runtime} & \multicolumn{2}{c}{Annotation contract}\\
Method & Paired & Precision & Paired & Precision\\
\midrule
Mypy & xxx & xxx & xxx & xxx\\
Pyright & xxx & xxx & xxx & xxx\\
Pyrefly & xxx & xxx & xxx & xxx\\
Pyinder & xxx & xxx & xxx & xxx\\
MOPSA (subset) & xxx & xxx & xxx & xxx\\
Direct LLM & xxx & xxx & xxx & xxx\\
Inference + checker & xxx & xxx & xxx & xxx\\
Our approach & xxx & xxx & xxx & xxx\\
\bottomrule
\end{tabular}
\end{table}

Table~\ref{tab:resolution} places detection beside resolution and cost.
Our analysis completes xxx\% of revision analyses and resolves xxx\%
of eligible targets. The unresolved, timed-out, and unsupported counts
are xxx, xxx, and xxx. The median and 95th-percentile times are xxx and
xxx seconds, and peak memory is xxx. Across all runs, the total number
of LLM requests is xxx and token consumption is xxx.

\begin{table}[t]
\centering
\small
\setlength{\tabcolsep}{3pt}
\caption{Resolution and cost. U, T, and S denote unresolved, timed-out, and unsupported targets. Time is the median per revision in seconds; total cost includes all runs. Outcome counts follow the aggregation rule in Section~\ref{sec:setup}.}
\label{tab:resolution}
\begin{tabular}{lrrrrrr}
\toprule
Method & Resolved (\%) & U & T & S & Time & Total (h)\\
\midrule
Mypy & xxx & xxx & xxx & xxx & xxx & xxx\\
Pyright & xxx & xxx & xxx & xxx & xxx & xxx\\
Pyrefly & xxx & xxx & xxx & xxx & xxx & xxx\\
Pyinder & xxx & xxx & xxx & xxx & xxx & xxx\\
MOPSA (subset) & xxx & xxx & xxx & xxx & xxx & xxx\\
Direct LLM & xxx & xxx & xxx & xxx & xxx & xxx\\
Inference + checker & xxx & xxx & xxx & xxx & xxx & xxx\\
Our approach & xxx & xxx & xxx & xxx & xxx & xxx\\
\bottomrule
\end{tabular}
\end{table}

We associate targets with the information needed to establish the checked
state: input/return dependencies, condition--type relations, and attribute
or global updates. These labels can overlap. For each differing outcome,
we trace the operands and effects supporting the judgment. A difference
is attributed to relation preservation when restoring the relation in a
matched mechanism comparison changes the decision. The counts for these
three dependence categories are xxx, xxx, and xxx; their respective paired
detections are xxx, xxx, and xxx.

\subsection{RQ2: Mechanism Effects and Cost}
\label{sec:rq2}
We compare controlled alternatives under the same entries, available
repository facts, operation semantics, model configuration, and budgets.
All variants use the same confirmation rules. Table~\ref{tab:variants}
defines the interventions; Table~\ref{tab:mechanism-results} reports
category-specific detections, unresolved targets, and total time.

\begin{table}[t]
\centering
\small
\caption{Controlled mechanism comparisons. Each variant changes the named representation or procedure while retaining the common reporting criteria.}
\label{tab:variants}
\begin{tabularx}{\linewidth}{@{}lX@{}}
\toprule
Variant & Change\\
\midrule
Independent types & Project each case's location relations onto independent type sets.\\
Merged paths & Join branch alternatives, removing condition--effect associations.\\
Unknown effects & Replace potentially written shared-state contents by unknown values.\\
Full-body construction & Construct relations from the whole function without separate entry-independent fact inference.\\
No implicit alternatives & Disable additional LLM type-combination proposals while retaining symbolic requirements and repository facts.\\
Relational summaries & Export intermediate states, normal/failure relations, and error locations; compose by linking intermediate symbols.\\
Direct body analysis & Analyze callee bodies at calls using the same state and operation rules.\\
No instance cache & Recompute instantiated results while retaining reusable summary templates.\\
Full callee expansion & Re-infer complete relevant callee bodies instead of selectively expanded type paths.\\
No dependency feedback & Assess initial candidates without revising summaries and reanalyzing dependent callers.\\
\bottomrule
\end{tabularx}
\end{table}

\begin{table}[t]
\centering
\small
\caption{Mechanism results. Runtime and contract columns report paired detections for their respective error categories. Outcome counts follow the aggregation rule in Section~\ref{sec:setup}; total cost includes all runs.}
\label{tab:mechanism-results}
\begin{tabular}{lrrrr}
\toprule
Variant & Runtime & Contract & Unresolved & Total (h)\\
\midrule
Complete approach & xxx & xxx & xxx & xxx\\
Independent types & xxx & xxx & xxx & xxx\\
Merged paths & xxx & xxx & xxx & xxx\\
Unknown effects & xxx & xxx & xxx & xxx\\
Full-body construction & xxx & xxx & xxx & xxx\\
No implicit alternatives & xxx & xxx & xxx & xxx\\
Relational summaries & xxx & xxx & xxx & xxx\\
Direct body analysis & xxx & xxx & xxx & xxx\\
No instance cache & xxx & xxx & xxx & xxx\\
Full callee expansion & xxx & xxx & xxx & xxx\\
No dependency feedback & xxx & xxx & xxx & xxx\\
\bottomrule
\end{tabular}
\end{table}

\paragraph{State and condition preservation.}
Independent types, merged paths, and unknown effects isolate operand
relations, branch correlations, and shared-state updates. Relative to the complete approach, their differences in paired runtime
detection counts are xxx, xxx, and xxx; the corresponding differences in
paired annotation-contract detection counts are xxx, xxx, and xxx. They
change xxx, xxx, and xxx unresolved outcomes, respectively.
Decision-level inspection links each change to the condition, value
version, or object binding that reaches the target operation.

\paragraph{Construction and reuse.}
Full-body construction retains the same type candidates and source
semantics, making the cost comparison sensitive to the reuse of
entry-independent facts. Shared construction produces xxx use-site facts
and consumes xxx seconds, compared with xxx seconds for full-body
construction. Disabling implicit alternatives changes the number of
represented type combinations by xxx. The corresponding differences in
paired runtime and annotation-contract detection counts are xxx and xxx.
The instance cache serves xxx\% of eligible applications. Direct body
analysis is evaluated with and without instance caching; the corresponding
total costs are xxx and xxx hours.

\paragraph{Ordered relations and refinement.}
The relational-summary comparison preserves intermediate failures and
error locations, so it compares how check dependencies are represented
and used. It produces xxx differing target decisions and xxx differing
unresolved outcomes. The selective and full-callee expansion variants use
the same re-inference procedure and contextual information. They expand
xxx and xxx source statements per candidate and consume xxx and xxx
seconds in conflict analysis. Dependency feedback changes xxx candidates
from unresolved to decided, dismisses xxx candidates after relation
correction, and establishes xxx runtime and xxx annotation-contract
reports. These outcomes are recorded with the responsible summary revision
and the affected caller judgment.

\subsection{RQ3: New-Error Discovery}
\label{sec:rq3}
The discovery corpus contains xxx repositories independent of the
historical benchmark, at revisions xxx. Their selection criteria and
analysis entries are xxx. Reports are reviewed using the same category
and evidence rules as the historical study. Across these repositories,
our approach identifies xxx runtime type errors and xxx
annotation-contract errors. Their diagnostic precision is xxx\% and
xxx\%, respectively. The numbers with independent maintainer confirmation
are xxx and xxx; the numbers without that confirmation are xxx and xxx.

For each confirmed error, the evidence record associates the source
location with the entry state, relevant calls, conditional updates, and
conflicting requirement or declaration. The discovery results contain
xxx errors involving attribute or global updates, xxx involving
condition--type relations, and xxx involving input/return dependencies.
Overlapping mechanism labels remain linked to the same error record.
Total analysis time over all runs is xxx hours, including xxx hours
of construction and xxx hours of conflict analysis and reanalysis.
